\cotefile{29}
\scanpage{188}
Opérateur d'\uline{homotopie} pour $d''$

% The pencilled wave cancels the two upper lines. A large X crosses the draft below.
\[
\begin{aligned}
&\deleted{\overline W_n(G)=G_{n-1}\times G_{n-2}\times\cdots\times G_0}\\
&\deleted{\deleted{\unsure{\overline W_n(G)|W_n(G)}}\ \tau:\overline W_n(G)\longrightarrow G_{n-1}}
\end{aligned}
\]

\[
\begin{aligned}
&\deleted{\varphi_n(x_1,\ldots,x_n)=\frac1n\left(1+\frac{p_1}{n-1}
+\frac{p_2}{\frac{(n-1)(n-2)}2}\,+\cdots+
\frac{p_k}{\binom{n-1}{k}}+\cdots+p_{n-1}\right)}\\
&\deleted{\text{où}\quad p_k=\deleted{\text{somme des}}\ \sum x_1 x_2\ldots x_k.}
\end{aligned}
\]

\[
\begin{aligned}
&\deleted{U_j f=\frac{1}{2\pi i}\int_{B_j} f(z_1,\ldots,z_{j-1},t,z_{j+1},\ldots,z_n)\,\frac{dt}{t-z_j}}\\
&\deleted{T_j f=\frac{1}{2\pi i}\iint_{K_j} f(z_1,\ldots,z_{j-1},t,z_{j+1},\ldots,z_n)\,\frac{dt\,d\bar t}{t-z_j}}
\end{aligned}
\]

\deleted{\deleted{\unsure{On|Et}} Pour \deleted{toute}toute suite strictement croissante $j_1<j_2<\ldots<j_q$, on \deleted{\unsure{note}} \inserted{introduit l'opérateur}
$\varphi_n(U\,;j_1,\ldots,j_q)\deleted{=\unsure{\psi_n(}}\inserted{=}$ ce que devient $\psi_n(U_1,\ldots,U_n)$ quand on remplace $U_{j_1}$, \ldots et $U_{j_q}$ par \uline{zéro}.}

\deleted{L'opérateur d'homotopie, pour $\omega=a\,d\bar z_{j_1}\wedge\cdots\wedge d\bar z_{j_q}$ \quad $(q\geq1)$ est}
\margin{\deleted{($n$ est le nombre des variables)}}
\[
\deleted{\underline{\left|\vphantom{\sum^q}\right. P\omega=\varphi_n(U\,;j_1,\ldots,j_q)\,\deleted{\unsure{\sum|\ast}}\sum_{\unsure{\nu|v|\gamma}=1}^{q}(-1)^{\nu-1}
(T_{j_\nu}a)\,d\bar z_{j_1}\wedge\cdots\wedge\widehat{d\bar z_{j_\nu}}\wedge\cdots\wedge d\bar z_{j_q}}}
\]

\[
\deleted{P d'' f\ \deleted{d''P f}=f-U_1 U_2\ldots U_n f}
\]
% archive sheet number: 187
